\section{Related Work}
\label{sec:related-ra}

The construction is bounded by three literatures: transaction
processing and the saga pattern, capability revocation in operating
systems, and endpoint detection and response engineering.

\paragraph{Compensating transactions and sagas.}
The saga pattern, introduced for long-running transactions whose
commit boundary is weaker than serializable schedules, pairs each
forward step with a compensating step that semantically undoes the
forward step~\cite{TODO_bea_bib}. Workflow nets extend the model with
timed arcs and soundness criteria deciding whether every reachable
marking has a path to a final marking. The construction here inherits
the saga vocabulary on the four reversible action variants: the
inverse executor is the compensating step, and the rollback receipt
is the typed artifact recording its completion. The construction does
not inherit the workflow-net machinery; the substrate has no timed
transition system beyond the monotonic counter, and the soundness
criterion the workflow-net literature proves is a strict generalization
the substrate does not claim. The contribution relative to sagas is
the constitutional-admission layer above the compensating step: a
rollback receipt is admissible only when the polity's predicates accept
its canonical body, which the saga literature treats as a workflow-
engine concern rather than as a polity-level decision.

\paragraph{Capability revocation in operating systems.}
EROS, seL4, and Capsicum treat capability revocation as a typed
operation on a capability token. The capability becomes unusable after
revocation, and the revocation event is a checkable artifact at the
machine boundary. The construction here lifts the same discipline to
the cross-organizational boundary where the receipt graph crosses
kernels: a revoke-grant action is destructive, requires bilateral
cosignature, and the constitutional admission predicate gates the
revocation. The OS-capability literature solves revocation at one
kernel; the construction here gates revocation at the receipt-graph
boundary between kernels.

\paragraph{Endpoint detection and response.}
Vendor EDR engines have shipped TTL fields and manual rollback
workflows for a decade. CrowdStrike Falcon, SentinelOne, Microsoft
Defender for Endpoint, and Sublime Security each record an action
identifier, target host, action class, TTL duration, and rollback
handle for response actions; auditors inspect the rollback handle by
hand when something needs to be reverted~\cite{TODO_bea_bib}. The
typed four-receipt grammar (request, execution, rollback,
acknowledgement) is a normalization of existing vendor schemas; the
constitutional admission predicate on the rollback receipt is the
substrate-level addition. The construction does not propose a new EDR;
it proposes a typed substrate the existing engines could in principle
adopt.

\paragraph{Provenance and lineage.}
Why-provenance, where-provenance, and how-provenance distinctions in
the database literature~\cite{TODO_bea_bib} are structurally related to
the receipt graph's lineage relations. A rollback receipt is a how-
provenance fragment: it records the inverse executor's signed
assertion that the OS operation completed, and the receipt graph
records the lineage from request to execution to rollback. The
construction does not extend the provenance literature; it specializes
it to the cross-kernel-admission case the parent paper already named.

\paragraph{Verified systems.}
The Lean-based industrial verification pattern is established by AWS
Cedar (Lean specification plus differential-tested Rust runtime) and
SampCert (Lean-verified differential-privacy primitives deployed
inside AWS Clean Rooms)~\cite{TODO_bea_bib}. The parent paper's
treaty intersection and amendment refinement theorems extend the
pattern to cross-kernel admission. The construction here adds the
TTL-bounded amendment chain theorem to the substrate, in the same
shape as the parent's essential-predicate chain.

\paragraph{Agentic-AI safety.}
The alignment-evaluation literature (METR, UK AI Safety Institute
Inspect, ARC Evals)~\cite{TODO_bea_bib} bounds what a model is likely
to do during pre-deployment evaluation; the construction here gates
whether a deployed agent's destructive action crosses the trust
boundary at runtime. The constructions compose rather than compete:
an evaluation result is an input predicate to the polity's
constitution, and the constitution's response to a destructive action
includes the bilateral-cosignature reduction. Constitutional AI
conditions outputs at training; the substrate here gates outputs at
deployment.
